\chapter{Conclusion}

% Understanding the evolutionary dynamics of Plasmodium and its vectors is critical for effective vector control

% spatial pop gen can help inform this

In this dissertation, I use spatial patterns of variation in  
\textit{Anopheles} malaria vectors in order to 
uncover spatial evolutionary dynamics that drive adaptation and divergence across their ranges.

In Chapter II, I present our work using correlated patterns of machine learning
prediction errors to infer common patterns of dispersal between \textit{Anopheles gambiae} 
mosquitoes and their \textit{Plasmodium falciparum} parasites. 
Being concerned about how the spatial distribution of training data impacted 
and biased our inferences, I assessed and then developed methods to reduce
these effects. When correcting for spatial bias in empirical sampling data,
correlations between \textit{Anopheles gambiae} and \textit{Plasmodium falciparum} residuals persist,
suggesting that these prediction errors reflect shared dispersal patterns 
between host and parasite.a

Given these findings, I became interested in how adaptive alleles
are spreading through these spatially-structured populations. In
Chapter III, I discover a spatial signal of positive selection through 
continuous-space simulation and use this observation to develop a method for 
detecting adaptive alleles in spatially-structured populations.
Applying this method to \textit{Anopheles gambiae}, I identify several 
novel alleles that appear to be under positive selection, including several 
at loci associated with insecticide resistance.